%! Characteristics of Linear Functions — Unit 1, Lesson 1.0 — Slides (Beamer)
% Deck follows the lesson's gradual-release flow: hook → warm-up → guided notes (I do / we do) →
% group activity → debrief → close & homework.
\documentclass[aspectratio=169, 11pt]{beamer}
\usepackage{algebra2-beamer}   % provides \forestheader{} and \sectionlabel[color]{}
% Standard coordinate grid + axes: {xmin}{xmax}{ymin}{ymax}
\newcommand{\lgrid}[4]{%
  \draw[step=1, linegray!45, thin] (#1,#3) grid (#2,#4);
  \draw[->, thick, charcoal] (#1,0)--(#2,0) node[below right, font=\tiny]{$x$};
  \draw[->, thick, charcoal] (0,#3)--(0,#4) node[left, font=\tiny]{$y$};}
% Temperature grid used by the hook and the activity: h on 0..7, T on -8..8.
\newcommand{\tempgrid}{%
  \draw[step=1, linegray!45, thin] (0,-8) grid (7,8);
  \draw[->, thick, charcoal] (0,0)--(7.6,0) node[right, font=\tiny]{$h$};
  \draw[->, thick, charcoal] (0,-8.6)--(0,8.8) node[above, font=\tiny]{$T$};
  \foreach \x in {1,...,7} \node[charcoal, font=\tiny, below] at (\x,0){\x};
  \foreach \y in {-8,-4,4,8} \node[charcoal, font=\tiny, left] at (0,\y){\y};}

\begin{document}

% TITLE SLIDE (CourseName is not defined in beamer, so the name is literal)
{
\setbeamercolor{background canvas}{bg=forest}
\begin{frame}[plain]
  \vspace{2.8cm}\hspace{0.55cm}%
  \begin{minipage}{0.9\textwidth}
    {\color{goldacc}\bfseries\small LESSON 1.0}\\[0.5cm]
    {\color{white}\bfseries\LARGE Characteristics of Linear Functions}\\[1.2cm]
    {\color{forestmid}\small Algebra 2~$\cdot$~Unit 1: Linear Functions}
  \end{minipage}
\end{frame}
}

% LEARNING TARGETS
\begin{frame}[t, plain]
  \forestheader{Today's targets}
  \sectionlabel{What you will be able to do}
  \begin{itemize}\setlength{\itemsep}{6pt}
    \item Find the \textbf{slope} of a line --- its constant \textbf{rate of change} --- from a
          graph, a table, or any two points.
    \item Find the \textbf{$y$-intercept} and the \textbf{zero}, and say what each one means in a
          situation.
    \item Say where a line is \textbf{increasing / decreasing / constant}, and where it is
          \textbf{positive / negative}.
    \item State the \textbf{domain} and \textbf{range} --- for the rule, and for the story.
  \end{itemize}
  \vspace{4pt}
  \begin{block}{How today runs}
    Warm-up (5) $\to$ notes together (15) $\to$ group activity (25) $\to$ debrief (10) $\to$ close (5).
  \end{block}
\end{frame}

% WARM-UP
\begin{frame}[t, plain]
  \forestheader{Warm-Up: three things you already know}
  \sectionlabel{5 minutes --- alone, then check with a neighbor}
  \begin{columns}[T]
    \begin{column}{0.55\textwidth}
      \begin{enumerate}\setlength{\itemsep}{6pt}
        \item Read the coordinates of $A$, $B$, $C$, $D$. Which one is \emph{on} the $y$-axis?
              On the $x$-axis?
        \item $g(x)=3x-6$: find $g(0)$, $g(2)$, $g(4)$ and write each as a point. Then: for what
              $x$ is $g(x)=0$?
        \item Continue the table $5, 8, 11, \dots$ and write a rule for $y$.
      \end{enumerate}
    \end{column}
    \begin{column}{0.42\textwidth}
      \centering
      \begin{tikzpicture}[scale=0.42]
        \lgrid{-4}{4}{-4}{4}
        \fill[forest] (3,2) circle (4pt)   node[above right, font=\scriptsize, charcoal]{$A$};
        \fill[forest] (-2,-3) circle (4pt) node[below left, font=\scriptsize, charcoal]{$B$};
        \fill[forest] (0,-1) circle (4pt)  node[below right, font=\scriptsize, charcoal]{$C$};
        \fill[forest] (-3,0) circle (4pt)  node[above left, font=\scriptsize, charcoal]{$D$};
      \end{tikzpicture}
    \end{column}
  \end{columns}
\end{frame}

% HOOK
\begin{frame}[t, plain]
  \forestheader{Hook: a cold Saturday morning}
  \sectionlabel[goldacc]{Two numbers tell the whole story}
  \begin{columns}[T]
    \begin{column}{0.54\textwidth}
      At 6:00 a.m.\ it was $-6\,^\circ$C, rising a steady $2\,^\circ$C an hour.
      \begin{itemize}\setlength{\itemsep}{4pt}
        \item $T(0)=-6$, \; $T(1)=-4$, \; $T(2)=-2$, \; $T(4)=2$
        \item It reaches $0\,^\circ$ at $h=3$ --- 9:00 a.m.
      \end{itemize}
      \vspace{3pt}
      \emph{Where it starts} and \emph{how fast it changes}. And one moment splits the morning into
      ``icy'' and ``not.''
    \end{column}
    \begin{column}{0.44\textwidth}
      \centering
      \begin{tikzpicture}[x=0.52cm, y=0.20cm]
        \tempgrid
        \draw[very thick, forest] (0,-6)--(7,8);
        \fill[charcoal] (0,-6) circle (3pt);
        \fill[charcoal] (3,0) circle (3pt);
      \end{tikzpicture}
    \end{column}
  \end{columns}
\end{frame}

% NOTES 1 — LINEAR / SLOPE
\begin{frame}[t, plain]
  \forestheader{Notes 1: linear means the rate of change is constant}
  \sectionlabel{Slope}
  \begin{columns}[T]
    \begin{column}{0.55\textwidth}
      The differences in the Saturday table are $+2$, $+2$, $+2$ --- \emph{constant}. That is what
      makes it \textbf{linear}, and the constant is the \textbf{slope}.
      \[ m=\frac{\Delta y}{\Delta x}=\frac{\text{rise}}{\text{run}}=\frac{y_2-y_1}{x_2-x_1}
         \qquad f(x)=mx+b \]
      From $(1,-2)$ and $(3,2)$: \; $m=\dfrac{2-(-2)}{3-1}=\dfrac{4}{2}=2$.
      \vspace{2pt}\\
      \emph{Any} two points work --- the rate is constant.
    \end{column}
    \begin{column}{0.42\textwidth}
      \centering
      \begin{tikzpicture}[scale=0.40]
        \lgrid{-2}{6}{-6}{6}
        \draw[very thick, forest] (-1,-6)--(5,6);
        \draw[thick, navy, dashed] (2,0)--(3,0)--(3,2);
        \node[navy, font=\tiny] at (2.5,-0.8){run $1$};
        \node[navy, font=\tiny, right] at (3.1,1){rise $2$};
        \fill[charcoal] (0,-4) circle (4pt);
        \fill[charcoal] (2,0) circle (4pt);
      \end{tikzpicture}\\[2pt]
      {\scriptsize $f(x)=2x-4$}
    \end{column}
  \end{columns}
\end{frame}

% NOTES 2 — INTERCEPTS
\begin{frame}[t, plain]
  \forestheader{Notes 2: where the graph meets the axes}
  \sectionlabel{Intercepts and zeros}
  \begin{itemize}\setlength{\itemsep}{6pt}
    \item \textbf{$y$-intercept} $(0,b)$: set $x=0$. \; The \emph{starting value}.
    \item \textbf{$x$-intercept}, the \textbf{zero}, $(a,0)$: solve $f(x)=0$. \; The moment the
          quantity runs out --- or crosses over.
  \end{itemize}
  \vspace{4pt}
  \begin{block}{Worked: $f(x)=2x-4$}
    $y$-intercept $(0,-4)$, straight off the rule. \; Zero: $2x-4=0 \Rightarrow 2x=4 \Rightarrow x=2$,
    so $(2,0)$. Check both against the graph.
  \end{block}
\end{frame}

% NOTES 3 — MOVES vs SITS
\begin{frame}[t, plain]
  \forestheader{Notes 3: which way it \emph{moves}, where it \emph{sits}}
  \sectionlabel[redacc]{Two different questions}
  \begin{columns}[T]
    \begin{column}{0.48\textwidth}
      \textbf{Which way it moves}
      \begin{itemize}\setlength{\itemsep}{3pt}
        \item $m>0$: \textbf{increasing}
        \item $m<0$: \textbf{decreasing}
        \item $m=0$: \textbf{constant}
      \end{itemize}
    \end{column}
    \begin{column}{0.48\textwidth}
      \textbf{Where it sits}
      \begin{itemize}\setlength{\itemsep}{3pt}
        \item $f(x)>0$: \textbf{above} the $x$-axis
        \item $f(x)<0$: \textbf{below} it
        \item the sign switches only at a \textbf{zero}
      \end{itemize}
    \end{column}
  \end{columns}
  \vspace{6pt}
  \begin{block}{Careful: ``positive'' is not ``increasing''}
    Sunday: $g(h)=6-2h$. At $h=1$, $g(1)=4$ --- a \emph{positive} number --- and the temperature is
    \emph{falling}. Both at once.
  \end{block}
\end{frame}

% NOTES 4 — DOMAIN / RANGE + PRACTICE
\begin{frame}[t, plain]
  \forestheader{Notes 4: domain and range --- the math, and the story}
  \sectionlabel{And your turn}
  \begin{itemize}\setlength{\itemsep}{4pt}
    \item Not horizontal: domain \textbf{all reals}, range \textbf{all reals}.
    \item Horizontal $c(x)=b$: slope $0$, range $\{b\}$, and \emph{no} zero unless $b=0$ --- and then
          \emph{every} input is one.
    \item In a story, the domain shrinks to the inputs that make sense.
  \end{itemize}
  \vspace{5pt}
  \begin{block}{Your turn (practice box): $h(x)=-2x+6$}
    Slope? \; $y$-intercept? \; Increasing or decreasing? \; Solve $-2x+6=0$ for the zero. \;
    Where is $h$ positive? Negative? \; Domain and range?
  \end{block}
\end{frame}

% GROUP ACTIVITY LAUNCH
\begin{frame}[t, plain]
  \forestheader{Group Activity: \emph{Freezing Point}}
  \sectionlabel[goldacc]{Groups of 3--4 $\cdot$ 25 minutes}
  \begin{columns}[T]
    \begin{column}{0.56\textwidth}
      \begin{enumerate}\setlength{\itemsep}{4pt}
        \item \textbf{Saturday morning} --- rising $2^\circ$/hr from $-6^\circ$. Table, rule, both
              intercepts, slope, sign intervals.
        \item \textbf{Sunday evening} --- falling $2^\circ$/hr from $6^\circ$. Same toolkit, plus
              the story-vs-rule domain question.
        \item \textbf{Side by side} --- same zero, opposite signs.
        \item \textbf{Model it} --- the \$40 car-wash card, $B(w)=40-8w$.
      \end{enumerate}
    \end{column}
    \begin{column}{0.42\textwidth}
      \centering
      \begin{tikzpicture}[x=0.50cm, y=0.19cm]
        \tempgrid
        \draw[very thick, forest] (0,-6)--(7,8);
        \draw[very thick, navy] (0,6)--(7,-8);
        \fill[charcoal] (3,0) circle (3pt);
      \end{tikzpicture}
    \end{column}
  \end{columns}
  \vspace{2pt}
  \textbf{The rule:} for every answer, be ready to show me \emph{where you see it on the graph}.
\end{frame}

% DEBRIEF
\begin{frame}[t, plain]
  \forestheader{Debrief: the four things to take away}
  \sectionlabel[redacc]{10 minutes --- correct your own page as we go}
  \begin{enumerate}\setlength{\itemsep}{5pt}
    \item \textbf{Part 1} labelled completely: slope $2$, $y$-intercept $(0,-6)$, zero $(3,0)$,
          negative before 9~a.m., positive after.
    \item \textbf{2c:} $4^\circ$ and falling. \emph{Positive} = where it \textbf{sits};
          \emph{increasing} = which way it \textbf{moves}.
    \item \textbf{Part 3:} both lines cross at $h=3$ --- so a line changes sign \emph{only} at its
          zero, and therefore at most once.
    \item \textbf{4d:} a \$10 wash does not move the start; it steepens the line, so the card runs
          out \emph{sooner} ($w=4$).
  \end{enumerate}
\end{frame}

% CLOSE
\begin{frame}[t, plain]
  \forestheader{Close}
  \sectionlabel{What changed today}
  You walked in able to read points off a graph. You walk out able to name what a line's features
  \emph{mean} --- where it starts, where it crosses zero, how fast it changes, which way it heads,
  and where it sits above or below the axis.
  \vspace{5pt}
  \begin{block}{Homework}
    The \textbf{Homework} page in your packet --- due next class. Six problems plus an extension:
    a rule, two graphs sharing a zero, a table (\emph{watch the step in $x$}), the flat line, a pool
    model, and one multiple-choice item.
  \end{block}
  \vspace{4pt}
  \textbf{Next:} Lesson 1.1 --- using these features to \emph{write} the equation of a line.
\end{frame}

\end{document}
